\centerline{\bf \Huge Турбо тренировка VII}
\centerline{\bf \Huge }

\begin{flushright}
        \scriptsize{}
\end{flushright}

\problem{1}
Пусть $a_1, a_2, a_3, \ldots$ --- бесконечная возрастающая последовательность натуральных чисел, а $p_1, p_2, p_3, \ldots-$ последовательность простых чисел такая, что при каждом натуральном $n$ число $a_n$ делится на $p_n$. Оказалось, что при всех натуральных $n$ и $k$ верно равенство $a_n-a_k=p_n-p_k$. Докажите, что все числа $a_1, a_2, \ldots$ простые.

\problem{2}
Из цифр $1,2,3,4,5,6,7,8,9$ составлены девять (не обязательно различных) девятизначных чисел; каждая из цифр использована в каждом числе ровно один раз. На какое наибольшее количество нулей может оканчиваться сумма этих девяти чисел?

\problem{3}
Окружность $\omega$ касается сторон $A B$ и $A C$ треугольника $A B C$. Окружность $\Omega$ касается стороны $A C$ и продолжения стороны $A B$ за точку $B$, а также касается $\omega$ в точке $L$, лежащей на стороне $B C$. Прямая $A L$ вторично пересекает $\omega$ и $\Omega$ в точках $K$ и $M$ соответственно. Оказалось, что $K B \| C M$. Докажите, что треугольник $L C M$ равнобедренный.

\problem{4}
Квадрат разбит на $n^2 \geqslant 4$ прямоугольников $2(n-1)$ прямыми, из которых $n-1$ параллельны одной стороне квадрата, а остальные $n-1$ - другой. Докажите, что можно выбрать $2 n$ прямоугольников разбиения таким образом, что для любых двух выбранных прямоугольников один из них можно поместить в другой (возможно, предварительно повернув).